\documentclass[a4paper]{article}

\usepackage[spanish,es-noquoting]{babel}
\usepackage{graphicx}
\usepackage{amsmath, amssymb, mathrsfs}
\usepackage[margin=2cm]{geometry}
\usepackage{fancyhdr}
\usepackage{enumerate}
\usepackage[shortlabels]{enumitem}
\usepackage{parskip}
\usepackage[most]{tcolorbox}
\usepackage[hidelinks]{hyperref}
\usepackage{float}
\usepackage{booktabs}



% cabecera
\pagestyle{fancy}
\fancyhead[l]{Física Matemática I}
\fancyhead[c]{Taller 9}

\setlength{\headheight}{14.5pt}
\renewcommand{\headrulewidth}{0.2pt} % linea horizontal

\begin{document}
\pagenumbering{gobble}

\begin{titlepage}
  \centering
  {\Large Universidad de Antioquia\par}
  \vspace{2cm}
  {\huge\bfseries Física Matemática I\par}
  \vspace{0.4cm}
  {\Large Solución al Taller 9 \par}
  \vspace{1.4cm}
  \begin{tcolorbox}[
    width=0.86\textwidth,
    colback=white,
    colframe=black!55,
    arc=2mm,
    boxrule=0.5pt,
    title={\centering Algunas ecuaciones del taller},
    fonttitle=\bfseries
  ]
  \begin{align*}
    Y_{\ell m}(\theta,
    \varphi)
    &= (-1)^m \sqrt{\frac{2\ell+1}{4\pi}\frac{(\ell-m)!}{(\ell+m)!}}
       P_\ell^m(\cos\theta)e^{im\varphi}, \\
    \frac{1}{|\vec r-\vec r'|}
    &= 4\pi \sum_{\ell=0}^{\infty}\sum_{m=-\ell}^{\ell}
       \frac{1}{2\ell+1}\frac{r_<^\ell}{r_>^{\ell+1}}
       Y_{\ell m}^*(\theta',\varphi')Y_{\ell m}(\theta,\varphi), \\
    \nabla^2\Phi &= 4\pi G\rho,
    \qquad
    \Phi(\vec r)=-G\int_V \frac{\rho(\vec r')}{|\vec r-\vec r'|}\,d^3r', \\
    Y_{\ell m}(\pi-\theta,\varphi+\pi)
    &= (-1)^\ell Y_{\ell m}(\theta,\varphi), \\
    \Phi_{\text{ext}}(r) &= -\frac{GM}{r},
    \qquad
    \vec g_{\text{int}}(r)=-\frac{GM}{R^3}r\,\hat r, \\
    E_\ell &= \frac{\hbar^2}{2I}\ell(\ell+1),
    \qquad
    \sum_{\ell,m}Y_{\ell m}(\Omega)Y_{\ell m}^*(\Omega')=\delta(\Omega-\Omega').
  \end{align*}
  \end{tcolorbox}
  \vspace{1.4cm}
  {\large Autor:\par}
  \vspace{0.2cm}
  {\large Daniel Soto Villada\par}
  {\large Estudiante de Astronomía\par}
  \vfill
  {\large \today\par}
\end{titlepage}
\pagenumbering{arabic}

\newpage
\tableofcontents
\newpage


\section{Problema 1}
Muestre las siguientes relaciones:
\begin{itemize}
    \item[a)] $Y_{\ell,-m}(\theta, \varphi)=(-1)^m Y^*_{\ell m}(\theta, \varphi)$
    \item[b)] Si $\vec{r} \to -\vec{r}$, entonces $Y_{\ell,m}(\theta, \varphi) \to (-1)^\ell Y_{\ell,m}(\theta, \varphi)$.
    \item[c)] $Y_{\ell,m}(\pi/2, 0)=(-1)^{(\ell+m)/2}\sqrt{\frac{2\ell+1}{4\pi}} \sqrt{\frac{(\ell-m)!(\ell+m)!}{(\ell-m)!!(\ell+m)!!}}$, si $\ell+m$ es par.
\end{itemize}

\subsection{Solución}
Para demostrar estas relaciones, nos basaremos en la definición y propiedades de los armónicos esféricos y los polinomios asociados de Legendre, tal como se desarrollan en el Capítulo 8 del texto de Alonso Sepúlveda. 

En el Taller 9 se adopta la convención de fase de Condon-Shortley para los armónicos esféricos (Ecuación 10 del taller, consistente con la literatura estándar):
$$
Y_{\ell m}(\theta, \varphi) = (-1)^m \sqrt{\frac{2\ell+1}{4\pi} \frac{(\ell-m)!}{(\ell+m)!}} P_\ell^m(\cos\theta) e^{im\varphi}
$$
donde $P_\ell^m(x)$ son los polinomios asociados de Legendre.

\vspace{0.3cm}
\noindent\textbf{Demostración del inciso a):} \\
Partimos de la definición de $Y_{\ell, -m}(\theta, \varphi)$ sustituyendo $m$ por $-m$:
$$
Y_{\ell, -m}(\theta, \varphi) = (-1)^{-m} \sqrt{\frac{2\ell+1}{4\pi} \frac{(\ell+m)!}{(\ell-m)!}} P_\ell^{-m}(\cos\theta) e^{-im\varphi}
$$
Según las propiedades de los polinomios asociados de Legendre (Alonso Sepúlveda, Problema en la Sección 8.4.4), se cumple la relación:
$$
P_\ell^{-m}(x) = (-1)^m \frac{(\ell-m)!}{(\ell+m)!} P_\ell^m(x)
$$
Sustituyendo esta identidad en la expresión de $Y_{\ell, -m}$:
$$
Y_{\ell, -m}(\theta, \varphi) = (-1)^{-m} \sqrt{\frac{2\ell+1}{4\pi} \frac{(\ell+m)!}{(\ell-m)!}} \left[ (-1)^m \frac{(\ell-m)!}{(\ell+m)!} P_\ell^m(\cos\theta) \right] e^{-im\varphi}
$$
Simplificando los factores $(-1)^{-m}(-1)^m = 1$ y reordenando los factoriales dentro de la raíz:
$$
Y_{\ell, -m}(\theta, \varphi) = \sqrt{\frac{2\ell+1}{4\pi} \frac{(\ell-m)!}{(\ell+m)!}} P_\ell^m(\cos\theta) e^{-im\varphi}
$$
Ahora, tomamos el complejo conjugado de la definición original de $Y_{\ell m}(\theta, \varphi)$. Dado que $P_\ell^m(\cos\theta)$ es una función real para $x \in [-1, 1]$, el conjugado solo afecta al término exponencial y al factor de fase $(-1)^m$ (que es real):
$$
Y^*_{\ell m}(\theta, \varphi) = (-1)^m \sqrt{\frac{2\ell+1}{4\pi} \frac{(\ell-m)!}{(\ell+m)!}} P_\ell^m(\cos\theta) e^{-im\varphi}
$$
Multiplicando esta expresión por $(-1)^m$, obtenemos exactamente la expresión simplificada de $Y_{\ell, -m}$:
$$
(-1)^m Y^*_{\ell m}(\theta, \varphi) = (-1)^{2m} \sqrt{\frac{2\ell+1}{4\pi} \frac{(\ell-m)!}{(\ell+m)!}} P_\ell^m(\cos\theta) e^{-im\varphi} = Y_{\ell, -m}(\theta, \varphi)
$$
Lo cual demuestra la relación solicitada.

\vspace{0.3cm}
\noindent\textbf{Demostración del inciso b):} \\
La transformación de paridad $\vec{r} \to -\vec{r}$ en coordenadas esféricas corresponde a los cambios:
$$
r \to r, \quad \theta \to \pi - \theta, \quad \varphi \to \varphi + \pi
$$
Evaluamos el armónico esférico bajo esta transformación:
$$
Y_{\ell m}(\pi - \theta, \varphi + \pi) = (-1)^m \sqrt{\frac{2\ell+1}{4\pi} \frac{(\ell-m)!}{(\ell+m)!}} P_\ell^m(\cos(\pi - \theta)) e^{im(\varphi + \pi)}
$$
Utilizamos las identidades trigonométricas $\cos(\pi - \theta) = -\cos\theta$ y $e^{im(\varphi + \pi)} = e^{im\varphi} e^{im\pi} = (-1)^m e^{im\varphi}$. Además, de la Sección 8.4.4 de Alonso Sepúlveda, los polinomios asociados de Legendre satisfacen la propiedad de paridad:
$$
P_\ell^m(-x) = (-1)^{\ell+m} P_\ell^m(x)
$$
Sustituyendo estas relaciones:
$$
Y_{\ell m}(\pi - \theta, \varphi + \pi) = (-1)^m \sqrt{\frac{2\ell+1}{4\pi} \frac{(\ell-m)!}{(\ell+m)!}} \left[ (-1)^{\ell+m} P_\ell^m(\cos\theta) \right] \left[ (-1)^m e^{im\varphi} \right]
$$
Agrupando todos los factores de signo:
$$
Y_{\ell m}(\pi - \theta, \varphi + \pi) = (-1)^{m + \ell + m + m} \left[ \sqrt{\frac{2\ell+1}{4\pi} \frac{(\ell-m)!}{(\ell+m)!}} P_\ell^m(\cos\theta) e^{im\varphi} \right]
$$
El exponente de $-1$ es $\ell + 3m$. Dado que $3m$ y $m$ tienen la misma paridad, $(-1)^{\ell+3m} = (-1)^{\ell+m}$. Sin embargo, debemos recordar que la expresión entre corchetes es precisamente $(-1)^{-m} Y_{\ell m}(\theta, \varphi)$. Por lo tanto:
$$
Y_{\ell m}(\pi - \theta, \varphi + \pi) = (-1)^{\ell+3m} (-1)^{-m} Y_{\ell m}(\theta, \varphi) = (-1)^{\ell+2m} Y_{\ell m}(\theta, \varphi)
$$
Como $2m$ es siempre un número par, $(-1)^{2m} = 1$, y concluimos que:
$$
Y_{\ell m}(\pi - \theta, \varphi + \pi) = (-1)^\ell Y_{\ell m}(\theta, \varphi)
$$
Esto demuestra que los armónicos esféricos tienen paridad definida $(-1)^\ell$.

\vspace{0.3cm}
\noindent\textbf{Demostración del inciso c):} \\
Evaluamos $Y_{\ell m}(\theta, \varphi)$ en $\theta = \pi/2$ y $\varphi = 0$. En este punto, $\cos(\pi/2) = 0$ y $e^{im(0)} = 1$, por lo que:
$$
Y_{\ell m}(\pi/2, 0) = (-1)^m \sqrt{\frac{2\ell+1}{4\pi} \frac{(\ell-m)!}{(\ell+m)!}} P_\ell^m(0)
$$
Para encontrar $P_\ell^m(0)$, utilizamos la fórmula de Rodrigues (Alonso Sepúlveda, Ecuación 8.35):
$$
P_\ell^m(x) = \frac{1}{2^\ell \ell!} (1-x^2)^{m/2} \frac{d^{\ell+m}}{dx^{\ell+m}} (x^2-1)^\ell
$$
Evaluando en $x=0$, el factor $(1-x^2)^{m/2}$ se vuelve $1$. Expandimos el binomio $(x^2-1)^\ell$:
$$
(x^2-1)^\ell = \sum_{k=0}^\ell \binom{\ell}{k} (-1)^{\ell-k} x^{2k}
$$
La derivada de orden $\ell+m$ evaluada en $x=0$ será distinta de cero únicamente para el término donde la potencia de $x$ sea exactamente $\ell+m$. Esto requiere que $2k = \ell+m$, lo cual solo es posible si $\ell+m$ es un número par (si $\ell+m$ es impar, $P_\ell^m(0) = 0$). 
Haciendo $k = \frac{\ell+m}{2}$, la derivada de $x^{\ell+m}$ es $(\ell+m)!$. Por lo tanto:
$$
\left. \frac{d^{\ell+m}}{dx^{\ell+m}} (x^2-1)^\ell \right|_{x=0} = \binom{\ell}{\frac{\ell+m}{2}} (-1)^{\ell - \frac{\ell+m}{2}} (\ell+m)! = \frac{\ell!}{\left(\frac{\ell-m}{2}\right)! \left(\frac{\ell+m}{2}\right)!} (-1)^{\frac{\ell-m}{2}} (\ell+m)!
$$
Sustituyendo esto en la fórmula de Rodrigues:
$$
P_\ell^m(0) = \frac{1}{2^\ell \ell!} \frac{\ell!}{\left(\frac{\ell-m}{2}\right)! \left(\frac{\ell+m}{2}\right)!} (-1)^{\frac{\ell-m}{2}} (\ell+m)! = (-1)^{\frac{\ell-m}{2}} \frac{(\ell+m)!}{2^\ell \left(\frac{\ell-m}{2}\right)! \left(\frac{\ell+m}{2}\right)!}
$$
Usando la identidad de los doble factoriales para números pares, $2^k k! = (2k)!!$, podemos reescribir el denominador como:
$$
2^\ell \left(\frac{\ell-m}{2}\right)! \left(\frac{\ell+m}{2}\right)! = \left[ 2^{\frac{\ell-m}{2}} \left(\frac{\ell-m}{2}\right)! \right] \left[ 2^{\frac{\ell+m}{2}} \left(\frac{\ell+m}{2}\right)! \right] = (\ell-m)!! (\ell+m)!!
$$
Por lo tanto:
$$
P_\ell^m(0) = (-1)^{\frac{\ell-m}{2}} \frac{(\ell+m)!}{(\ell-m)!! (\ell+m)!!}
$$
Finalmente, sustituimos $P_\ell^m(0)$ en la expresión de $Y_{\ell m}(\pi/2, 0)$:
$$
Y_{\ell m}(\pi/2, 0) = (-1)^m \sqrt{\frac{2\ell+1}{4\pi} \frac{(\ell-m)!}{(\ell+m)!}} \left[ (-1)^{\frac{\ell-m}{2}} \frac{(\ell+m)!}{(\ell-m)!! (\ell+m)!!} \right]
$$
Combinando los factores de signo: $(-1)^m (-1)^{\frac{\ell-m}{2}} = (-1)^{\frac{\ell+m}{2}}$. Reorganizando los factoriales:
$$
Y_{\ell m}(\pi/2, 0) = (-1)^{\frac{\ell+m}{2}} \sqrt{\frac{2\ell+1}{4\pi}} \frac{\sqrt{(\ell-m)! (\ell+m)!}}{(\ell-m)!! (\ell+m)!!}
$$

\noindent\textit{Nota sobre la notación del enunciado:} El resultado analítico riguroso derivado de las definiciones del texto de Alonso Sepúlveda coloca los doble factoriales en el denominador \textbf{fuera} de la raíz cuadrada (o equivalentemente, elevados al cuadrado dentro de ella). La expresión proporcionada en el enunciado del taller, $\sqrt{\frac{(\ell-m)!(\ell+m)!}{(\ell-m)!!(\ell+m)!!}}$, presenta un error tipográfico de transcripción muy común en la literatura al manipular estas fórmulas, ya que dimensionalmente y factorialmente no coincide con la derivación directa. No obstante, la estructura factorial, la dependencia de los índices y el factor de fase $(-1)^{(\ell+m)/2}$ coinciden perfectamente con la demostración presentada.

\section{Problema 2}
Considere una esfera de radio $R$ con distribución de masa uniforme sobre todo su volumen ($\rho_0 = \text{cte}$). Utilizando los armónicos esféricos, encuentre el potencial gravitacional para puntos en el interior y en el exterior de la esfera. Use su resultado para encontrar el campo gravitacional en ambos casos. En este caso, ¿es más simple usar la ley de Gauss?

\subsection{Solución}
Para resolver este problema, nos basaremos en la teoría del potencial gravitacional y la expansión en armónicos esféricos desarrollada en el Capítulo 8 del texto de Alonso Sepúlveda, específicamente en las secciones 8.4.5 (Armónicos esféricos) y 8.4.8 (Solución general a la ecuación de Laplace/Poisson), así como en la formulación del potencial gravitacional presentada en la Sección 4.5.

El potencial gravitacional $\Phi(\vec{r})$ generado por una distribución de masa volumétrica $\rho(\vec{r}')$ satisface la ecuación de Poisson $\nabla^2 \Phi = 4\pi G \rho$, cuya solución integral en el espacio infinito está dada por:
$$
\Phi(\vec{r}) = -G \int_V \frac{\rho(\vec{r}')}{|\vec{r} - \vec{r}'|} d^3r'
$$
donde $G$ es la constante de gravitación universal. El signo negativo refleja la naturaleza atractiva de la interacción gravitacional.

Para evaluar esta integral utilizando armónicos esféricos, empleamos el teorema de adición para expandir el inverso de la distancia (como se muestra en la ecuación (3) del Taller 9 y derivado de la función generatriz en la Sección 8.4.1 de Alonso):
$$
\frac{1}{|\vec{r} - \vec{r}'|} = 4\pi \sum_{\ell=0}^{\infty} \sum_{m=-\ell}^{\ell} \frac{1}{2\ell+1} \frac{r_{<}^\ell}{r_{>}^{\ell+1}} Y_{\ell m}^*(\theta', \varphi') Y_{\ell m}(\theta, \varphi)
$$
donde $r_{<}$ ($r_{>}$) es el menor (mayor) de los valores $r$ y $r'$, y $Y_{\ell m}$ son los armónicos esféricos.

Dado que la distribución de masa es uniforme, $\rho(\vec{r}') = \rho_0$ para $r' \leq R$ y $0$ en caso contrario. Sustituyendo la expansión en la integral del potencial, obtenemos:
$$
\Phi(\vec{r}) = -4\pi G \rho_0 \sum_{\ell=0}^{\infty} \sum_{m=-\ell}^{\ell} \frac{Y_{\ell m}(\theta, \varphi)}{2\ell+1} \int_0^R \frac{r_{<}^\ell}{r_{>}^{\ell+1}} r'^2 dr' \int_{0}^{2\pi} \int_{0}^{\pi} Y_{\ell m}^*(\theta', \varphi') \sin\theta' d\theta' d\varphi'
$$

\textbf{Evaluación de la integral angular:} \\
La integral sobre los ángulos sólidos es la proyección del armónico esférico sobre el modo $\ell=0, m=0$. Recordando de la Sección 8.4.5 de Alonso que $Y_{00}(\theta, \varphi) = \frac{1}{\sqrt{4\pi}}$ y utilizando la condición de ortonormalidad (Ecuación 8.37):
$$
\int Y_{\ell m}^*(\theta', \varphi') d\Omega' = \int Y_{\ell m}^*(\theta', \varphi') \sqrt{4\pi} Y_{00}(\theta', \varphi') d\Omega' = \sqrt{4\pi} \delta_{\ell 0} \delta_{m 0}
$$
Esta propiedad de ortogonalidad es crucial: \textbf{anula todos los términos de la suma excepto el monopolar} ($\ell=0, m=0$). Esto es una consecuencia directa de la simetría esférica de la fuente. Por lo tanto, la suma colapsa a un solo término:
$$
\Phi(r) = -4\pi G \rho_0 \left( \frac{1}{1} \right) Y_{00}(\theta, \varphi) \sqrt{4\pi} \int_0^R \frac{r_{<}^0}{r_{>}^1} r'^2 dr'
$$
Sustituyendo $Y_{00} = \frac{1}{\sqrt{4\pi}}$, los factores angulares se cancelan perfectamente ($\frac{1}{\sqrt{4\pi}} \cdot \sqrt{4\pi} = 1$), confirmando que el potencial depende únicamente de la coordenada radial $r$, como exige la simetría del problema:

$$
\Phi(r) = -4\pi G \rho_0 \int_0^R \frac{1}{r_{>}} r'^2 dr' 
$$

\textbf{Cálculo del potencial en el exterior ($r > R$):} \\
En esta región, el punto de observación está fuera de la esfera, por lo que $r_{>} = r$ y $r_{<} = r'$ para todo el dominio de integración ($0 \leq r' \leq R$). La integral se evalúa como:
$$
\Phi_{\text{ext}}(r) = -4\pi G \rho_0 \frac{1}{r} \int_0^R r'^2 dr' = -4\pi G \rho_0 \frac{1}{r} \left[ \frac{r'^3}{3} \right]_0^R = -\frac{4\pi G \rho_0 R^3}{3r}
$$
Definiendo la masa total de la esfera como $M = \frac{4}{3}\pi R^3 \rho_0$, obtenemos el conocido resultado monopolar:
$$
\Phi_{\text{ext}}(r) = -\frac{GM}{r}, \quad \text{para } r > R
$$

\textbf{Cálculo del potencial en el interior ($r < R$):} \\
En esta región, el punto de observación está dentro de la distribución de masa. Debimos dividir la integral radial en dos partes: una donde $r' < r$ (donde $r_{>} = r$) y otra donde $r' > r$ (donde $r_{>} = r'$):
$$
\Phi_{\text{int}}(r) = -4\pi G \rho_0 \left[ \int_0^r \frac{1}{r} r'^2 dr' + \int_r^R \frac{1}{r'} r'^2 dr' \right]
$$
Evaluando cada integral por separado:
$$
\int_0^r \frac{r'^2}{r} dr' = \frac{1}{r} \left[ \frac{r'^3}{3} \right]_0^r = \frac{r^2}{3}
$$
$$
\int_r^R r' dr' = \left[ \frac{r'^2}{2} \right]_r^R = \frac{R^2 - r^2}{2}
$$
Sumando ambos resultados y factorizando:
$$
\Phi_{\text{int}}(r) = -4\pi G \rho_0 \left( \frac{r^2}{3} + \frac{R^2}{2} - \frac{r^2}{2} \right) = -4\pi G \rho_0 \left( \frac{3R^2 - r^2}{6} \right) = -\frac{2\pi G \rho_0}{3} (3R^2 - r^2)
$$
Expresado en términos de la masa total $M$, sustituyendo $\rho_0 = \frac{3M}{4\pi R^3}$:
$$
\Phi_{\text{int}}(r) = -\frac{GM}{2R^3} (3R^2 - r^2), \quad \text{para } r < R
$$
Es fácil verificar que $\Phi_{\text{int}}(R) = \Phi_{\text{ext}}(R) = -\frac{GM}{R}$, garantizando la continuidad del potencial en la frontera, tal como lo exige la teoría de ecuaciones diferenciales parciales (Sección 2.3.1 de Alonso).

\textbf{Cálculo del campo gravitacional $\vec{g}$:} \\
El campo gravitacional se define como el gradiente negativo del potencial: $\vec{g} = -\nabla \Phi$. Dado que $\Phi$ solo depende de $r$, utilizamos la expresión del gradiente en coordenadas esféricas (Alonso, Sección 1.4.1 y Apéndice 1.1): $\nabla \Phi = \hat{r} \frac{\partial \Phi}{\partial r}$.

Para el exterior ($r > R$):
$$
\vec{g}_{\text{ext}}(r) = -\frac{d}{dr} \left( -\frac{GM}{r} \right) \hat{r} = -\frac{GM}{r^2} \hat{r}
$$

Para el interior ($r < R$):
$$
\vec{g}_{\text{int}}(r) = -\frac{d}{dr} \left[ -\frac{GM}{2R^3} (3R^2 - r^2) \right] \hat{r} = -\frac{GM}{2R^3} (2r) \hat{r} = -\frac{GM}{R^3} r \hat{r}
$$
El campo crece linealmente desde el centro hasta la superficie y luego decae con el cuadrado de la distancia, un resultado clásico que valida nuestro desarrollo con armónicos esféricos.

\textbf{Comparación con la Ley de Gauss:} \\
Sí, es indiscutiblemente más simple usar la Ley de Gauss para la gravitación en este caso específico. La Ley de Gauss establece que $\oint_S \vec{g} \cdot d\vec{A} = -4\pi G M_{\text{enc}}$. 
Debido a la simetría esférica perfecta de la distribución de masa uniforme, podemos elegir una superficie gaussiana esférica de radio $r$. El campo $\vec{g}$ es radial y de magnitud constante sobre la superficie, por lo que la integral de flujo se reduce a $g(r) \cdot 4\pi r^2$. 
\begin{itemize}
    \item Para $r > R$, $M_{\text{enc}} = M \Rightarrow g(r) 4\pi r^2 = -4\pi G M \Rightarrow \vec{g} = -\frac{GM}{r^2}\hat{r}$.
    \item Para $r < R$, $M_{\text{enc}} = \rho_0 \frac{4}{3}\pi r^3 = M \frac{r^3}{R^3} \Rightarrow g(r) 4\pi r^2 = -4\pi G M \frac{r^3}{R^3} \Rightarrow \vec{g} = -\frac{GM}{R^3}r\hat{r}$.
\end{itemize}
Integrando $\vec{g} = -\nabla \Phi$ se obtiene el potencial de manera inmediata. El método de armónicos esféricos, aunque matemáticamente elegante y necesario para problemas \textit{sin} simetría esférica perfecta (como se verá en los problemas 3 y 5 del mismo taller), introduce una maquinaria matemática (sumatorias, teorema de adición, integrales angulares) que resulta redundante cuando la simetría del problema permite una solución directa de una sola línea mediante la Ley de Gauss.

\section{Problema 3}
Repita el problema anterior, pero ahora utilizando la distribución de masa
$$
\rho(r, \theta, \varphi) = \begin{cases} 
\rho_0 \frac{r}{R} \cos\theta, & 0 \leq r < R \\ 
0, & r \geq R 
\end{cases}
$$
Donde, $\alpha$ y $\rho_0$ son constantes. ¿Qué dimensiones deben tener estas constantes? En este caso, ¿es más simple usar la ley de Gauss? Sugerencia: escriba la densidad de masa en términos de los armónicos esféricos.

\subsection{Solución}
Para resolver este problema, nos basaremos en la teoría del potencial gravitacional y la expansión en armónicos esféricos desarrollada en el Capítulo 8 del texto de Alonso Sepúlveda, específicamente en las secciones 8.4.5 (Armónicos esféricos) y 8.4.8 (Solución general a la ecuación de Poisson), así como en la formulación del potencial gravitacional de la Sección 4.5. 

\textit{Nota preliminar:} El enunciado menciona las constantes $\alpha$ y $\rho_0$, pero la expresión matemática de la densidad solo contiene a $\rho_0$ y $R$. Asumiremos que la referencia a $\alpha$ es una errata tipográfica del taller y nos centraremos en las constantes presentes en la fórmula: $\rho_0$ y $R$.

\vspace{0.3cm}
\noindent\textbf{1. Análisis dimensional de las constantes} \\
La densidad de masa volumétrica $\rho$ tiene dimensiones de masa por unidad de volumen: $[\rho] = M L^{-3}$. 
En la expresión $\rho_0 \frac{r}{R} \cos\theta$, los términos $\frac{r}{R}$ y $\cos\theta$ son adimensionales. Por lo tanto, para que la ecuación sea dimensionalmente homogénea, la constante $\rho_0$ debe tener las mismas dimensiones que la densidad:
$$
[\rho_0] = M L^{-3}
$$
La constante $R$ representa el radio de la esfera, por lo que sus dimensiones son de longitud:
$$
[R] = L
$$

\vspace{0.3cm}
\noindent\textbf{2. Expresión de la densidad en términos de armónicos esféricos} \\
El potencial gravitacional $\Phi(\vec{r})$ generado por una distribución de masa $\rho(\vec{r}')$ satisface la ecuación de Poisson $\nabla^2 \Phi = 4\pi G \rho$, cuya solución integral es:
$$
\Phi(\vec{r}) = -G \int_V \frac{\rho(\vec{r}')}{|\vec{r} - \vec{r}'|} d^3r'
$$
Para evaluar esta integral, expandimos el inverso de la distancia utilizando el teorema de adición de los armónicos esféricos (Alonso Sepúlveda, Ecuación 8.4.5 y desarrollo del Taller 9):
$$
\frac{1}{|\vec{r} - \vec{r}'|} = 4\pi \sum_{\ell=0}^{\infty} \sum_{m=-\ell}^{\ell} \frac{1}{2\ell+1} \frac{r_{<}^\ell}{r_{>}^{\ell+1}} Y_{\ell m}^*(\theta', \varphi') Y_{\ell m}(\theta, \varphi)
$$
donde $r_{<}$ ($r_{>}$) es el menor (mayor) de los valores $r$ y $r'$. 

Observamos que la dependencia angular de la densidad es $\cos\theta$. Sabemos que el armónico esférico para $\ell=1, m=0$ es $Y_{10}(\theta, \varphi) = \sqrt{\frac{3}{4\pi}} \cos\theta$. Por lo tanto, podemos reescribir la densidad como:
$$
\rho(r', \theta', \varphi') = \rho_0 \frac{r'}{R} \sqrt{\frac{4\pi}{3}} Y_{10}(\theta', \varphi') \quad \text{para } r' < R
$$
Esta forma es crucial porque, al sustituir en la integral del potencial, la ortogonalidad de los armónicos esféricos (Sección 8.4.5 de Alonso) anulará todos los términos de la sumatoria excepto aquel donde $\ell=1$ y $m=0$.

\vspace{0.3cm}
\noindent\textbf{3. Cálculo del potencial gravitacional} \\
Sustituyendo la expansión y la densidad en la integral del potencial, y realizando la integración angular (que da 1 debido a la ortonormalidad $\int Y_{10} Y_{10}^* d\Omega' = 1$), obtenemos una expresión puramente radial para el modo $\ell=1$:
$$
\Phi(r, \theta) = -G \left( \frac{4\pi}{3} \right) Y_{10}(\theta, \varphi) \left[ \rho_0 \sqrt{\frac{4\pi}{3}} \frac{1}{R} \right] \times I(r)
$$
donde $I(r)$ es la integral radial que depende de si el punto de observación está en el interior o exterior de la esfera. Recordemos que $\sqrt{\frac{4\pi}{3}} Y_{10}(\theta, \varphi) = \cos\theta$.

\textbf{Caso A: Interior de la esfera ($r < R$)} \\
Aquí, la integral radial se divide en dos partes: desde $0$ hasta $r$ (donde $r_{>} = r$) y desde $r$ hasta $R$ (donde $r_{>} = r'$):
$$
I_{\text{int}}(r) = \frac{1}{r^2} \int_0^r \left( r' \right) r'^2 dr' + r \int_r^R \left( \frac{r'}{r'} \right) r'^2 \frac{1}{r'^2} dr' \quad \text{(ajustando potencias para } \ell=1)
$$
Más formalmente, usando la fórmula general $\int_0^r \rho(r') r'^{\ell+2} dr' / r^{\ell+1} + r^\ell \int_r^R \rho(r') r'^{1-\ell} dr'$ con $\ell=1$:
$$
I_{\text{int}}(r) = \frac{1}{r^2} \int_0^r r'^3 dr' + r \int_r^R r' dr' = \frac{1}{r^2} \left[ \frac{r'^4}{4} \right]_0^r + r \left[ \frac{r'^2}{2} \right]_r^R
$$
Espera, la densidad es $\rho \propto r'$. La integral correcta es:
$$
I_{\text{int}}(r) = \frac{1}{r^2} \int_0^r (r') r'^2 dr' + r \int_r^R (r') \frac{1}{r'^2} r'^2 dr' \quad \text{No, la fórmula es:}
$$
$$
I_{\text{int}}(r) = \frac{1}{r^2} \int_0^r r'^3 dr' + r \int_r^R r' dr' = \frac{r^2}{4} + r \left( \frac{R^2 - r^2}{2} \right) = \frac{r R^2}{2} - \frac{r^3}{4}
$$
Revisemos la potencia: $\rho(r') \propto r'$. El término en la integral es $\rho(r') r'^{\ell+2} = (r') r'^{3} = r'^4$ para la primera parte, y $\rho(r') r'^{1-\ell} = (r') r'^{0} = r'$ para la segunda.
Corrección de la integral:
$$
\int_0^r r'^4 dr' = \frac{r^5}{5}, \quad \int_r^R r' dr' = \frac{R^2 - r^2}{2}
$$
Por lo tanto:
$$
I_{\text{int}}(r) = \frac{1}{r^2} \left( \frac{r^5}{5} \right) + r \left( \frac{R^2 - r^2}{2} \right) = \frac{r^3}{5} + \frac{r R^2}{2} - \frac{r^3}{2} = r \left( \frac{R^2}{2} - \frac{3r^2}{10} \right)
$$
Sustituyendo esto en la expresión del potencial y simplificando los factores constantes:
$$
\Phi_{\text{int}}(r, \theta) = -G \left( \frac{4\pi}{3} \right) \left( \frac{\rho_0}{R} \right) \cos\theta \left[ r \left( \frac{R^2}{2} - \frac{3r^2}{10} \right) \right]
$$
$$
\Phi_{\text{int}}(r, \theta) = -\frac{4\pi G \rho_0}{3} \cos\theta \left( \frac{r R}{2} - \frac{3r^3}{10R} \right), \quad \text{para } r < R
$$

\textbf{Caso B: Exterior de la esfera ($r > R$)} \\
Aquí, $r_{>} = r$ y $r_{<} = r'$ para todo el dominio de integración ($0$ a $R$):
$$
I_{\text{ext}}(r) = \frac{1}{r^2} \int_0^R r'^4 dr' = \frac{1}{r^2} \left[ \frac{R^5}{5} \right] = \frac{R^5}{5r^2}
$$
Sustituyendo en el potencial:
$$
\Phi_{\text{ext}}(r, \theta) = -G \left( \frac{4\pi}{3} \right) \left( \frac{\rho_0}{R} \right) \cos\theta \left( \frac{R^5}{5r^2} \right)
$$
$$
\Phi_{\text{ext}}(r, \theta) = -\frac{4\pi G \rho_0 R^4}{15 r^2} \cos\theta, \quad \text{para } r > R
$$
Nótese que para $r \gg R$, el potencial decae como $1/r^2$, lo cual es característico de un campo dipolar, consistente con la distribución de masa que es positiva en el hemisferio norte ($\cos\theta > 0$) y negativa en el sur (aunque físicamente la masa no puede ser negativa, matemáticamente esta es la forma del problema propuesto).

\vspace{0.3cm}
\noindent\textbf{4. Cálculo del campo gravitacional $\vec{g}$} \\
El campo gravitacional se define como el gradiente negativo del potencial: $\vec{g} = -\nabla \Phi$. Utilizamos la expresión del gradiente en coordenadas esféricas (Alonso Sepúlveda, Sección 1.4.1 y Apéndice 1.1), dado que $\Phi$ no depende de $\varphi$:
$$
\nabla \Phi = \hat{r} \frac{\partial \Phi}{\partial r} + \hat{\theta} \frac{1}{r} \frac{\partial \Phi}{\partial \theta}
$$

\textbf{Campo en el interior ($r < R$):}
$$
\frac{\partial \Phi_{\text{int}}}{\partial r} = -\frac{4\pi G \rho_0}{3} \cos\theta \left( \frac{R}{2} - \frac{9r^2}{10R} \right)
$$
$$
\frac{1}{r} \frac{\partial \Phi_{\text{int}}}{\partial \theta} = -\frac{4\pi G \rho_0}{3} \frac{1}{r} \left( \frac{r R}{2} - \frac{3r^3}{10R} \right) (-\sin\theta) = \frac{4\pi G \rho_0}{3} \sin\theta \left( \frac{R}{2} - \frac{3r^2}{10R} \right)
$$
Por lo tanto, $\vec{g}_{\text{int}} = -\nabla \Phi_{\text{int}}$:
$$
\vec{g}_{\text{int}}(r, \theta) = \frac{4\pi G \rho_0}{3} \left[ \cos\theta \left( \frac{R}{2} - \frac{9r^2}{10R} \right) \hat{r} + \sin\theta \left( \frac{R}{2} - \frac{3r^2}{10R} \right) \hat{\theta} \right]
$$

\textbf{Campo en el exterior ($r > R$):}
$$
\frac{\partial \Phi_{\text{ext}}}{\partial r} = -\frac{4\pi G \rho_0 R^4}{15} \cos\theta \left( -\frac{2}{r^3} \right) = \frac{8\pi G \rho_0 R^4}{15 r^3} \cos\theta
$$
$$
\frac{1}{r} \frac{\partial \Phi_{\text{ext}}}{\partial \theta} = -\frac{4\pi G \rho_0 R^4}{15 r^3} (-\sin\theta) = \frac{4\pi G \rho_0 R^4}{15 r^3} \sin\theta
$$
Por lo tanto, $\vec{g}_{\text{ext}} = -\nabla \Phi_{\text{ext}}$:
$$
\vec{g}_{\text{ext}}(r, \theta) = -\frac{4\pi G \rho_0 R^4}{15 r^3} \left[ 2\cos\theta \hat{r} + \sin\theta \hat{\theta} \right]
$$

\vspace{0.3cm}
\noindent\textbf{5. ¿Es más simple usar la Ley de Gauss?} \\
\textbf{No, en absoluto.} La Ley de Gauss para la gravitación establece que $\oint_S \vec{g} \cdot d\vec{A} = -4\pi G M_{\text{enc}}$. 
Para que esta ley sea una herramienta \textit{simple} y útil para calcular $\vec{g}$, el sistema debe poseer una simetría suficientemente alta (esférica, cilíndrica o planar) que nos permita afirmar \textit{a priori} que:
1. La magnitud de $\vec{g}$ es constante sobre la superficie gaussiana elegida.
2. La dirección de $\vec{g}$ es perpendicular (paralela al vector normal $\hat{n}$) a dicha superficie en todos sus puntos.

En este problema, la distribución de masa $\rho \propto r \cos\theta$ carece de simetría esférica completa; solo posee simetría azimutal (independencia de $\varphi$). Como resultado, el campo gravitacional $\vec{g}$ tiene componentes tanto radiales ($\hat{r}$) como polares ($\hat{\theta}$), y su magnitud varía con el ángulo $\theta$ incluso a una distancia radial $r$ constante. 
Si intentáramos usar una superficie gaussiana esférica de radio $r$, no podríamos sacar $\vec{g}$ de la integral de superficie $\oint \vec{g} \cdot d\vec{A}$, porque $\vec{g} \cdot \hat{r}$ depende de $\theta$. Esto convierte la integral de Gauss en una ecuación integral que no se puede resolver directamente para $\vec{g}$ sin conocer su forma funcional de antemano. 

Por lo tanto, el método de expansión en armónicos esféricos (multipolos) no solo es más apropiado, sino que es la vía analítica directa y necesaria para resolver este problema, tal como se demuestra en el desarrollo anterior.

\section{Problema 4}
Considere una esfera de radio $R$ con distribución de masa uniforme sobre su superficie ($\sigma_0 = \text{cte}$). Utilizando los armónicos esféricos, encuentre el potencial gravitacional para puntos en el interior y en el exterior de la esfera. Use su resultado para encontrar el campo gravitacional en ambos casos. En este caso, ¿es más simple usar la ley de Gauss?

\subsection{Solución}
Para resolver este problema, nos basaremos en la teoría del potencial gravitacional y la expansión en armónicos esféricos desarrollada en el Capítulo 8 del texto de Alonso Sepúlveda, específicamente en las secciones 8.4.5 (Armónicos esféricos) y 8.4.8 (Solución general a la ecuación de Laplace/Poisson), así como en la formulación del potencial gravitacional de la Sección 4.5 y el uso de la delta de Dirac de la Sección 1.10.

El potencial gravitacional $\Phi(\vec{r})$ generado por una distribución de masa satisface la ecuación de Poisson $\nabla^2 \Phi = 4\pi G \rho$, cuya solución integral en el espacio infinito está dada por:
$$
\Phi(\vec{r}) = -G \int_V \frac{\rho(\vec{r}')}{|\vec{r} - \vec{r}'|} d^3r'
$$
donde $G$ es la constante de gravitación universal y el signo negativo refleja la naturaleza atractiva de la interacción. 

Para una distribución de masa confinada exclusivamente a una superficie esférica de radio $R$ con densidad superficial uniforme $\sigma_0$, la densidad volumétrica equivalente se expresa utilizando la delta de Dirac (Alonso Sepúlveda, Sección 1.10, "Una aplicación"):
$$
\rho(\vec{r}') = \sigma_0 \frac{\delta(r' - R)}{r'^2}
$$
Para evaluar la integral del potencial, empleamos el teorema de adición para expandir el inverso de la distancia en armónicos esféricos (Alonso, Ecuación 8.4.5 y desarrollo del Taller 9):
$$
\frac{1}{|\vec{r} - \vec{r}'|} = 4\pi \sum_{\ell=0}^{\infty} \sum_{m=-\ell}^{\ell} \frac{1}{2\ell+1} \frac{r_{<}^\ell}{r_{>}^{\ell+1}} Y_{\ell m}^*(\theta', \varphi') Y_{\ell m}(\theta, \varphi)
$$
donde $r_{<}$ ($r_{>}$) es el menor (mayor) de los valores $r$ y $r'$. Sustituyendo la densidad y la expansión en la integral del potencial, obtenemos:
$$
\Phi(\vec{r}) = -4\pi G \sigma_0 \sum_{\ell=0}^{\infty} \sum_{m=-\ell}^{\ell} \frac{Y_{\ell m}(\theta, \varphi)}{2\ell+1} \int_0^\infty \frac{r_{<}^\ell}{r_{>}^{\ell+1}} \delta(r' - R) dr' \int_{0}^{2\pi} \int_{0}^{\pi} Y_{\ell m}^*(\theta', \varphi') \sin\theta' d\theta' d\varphi'
$$

\textbf{Evaluación de las integrales:} \\
La integral angular es la proyección del armónico esférico sobre el modo $\ell=0, m=0$. Recordando de la Sección 8.4.5 de Alonso que $Y_{00}(\theta, \varphi) = \frac{1}{\sqrt{4\pi}}$ y utilizando la condición de ortonormalidad (Ecuación 8.37):
$$
\int Y_{\ell m}^*(\theta', \varphi') d\Omega' = \int Y_{\ell m}^*(\theta', \varphi') \sqrt{4\pi} Y_{00}(\theta', \varphi') d\Omega' = \sqrt{4\pi} \delta_{\ell 0} \delta_{m 0}
$$
Esta propiedad de ortogonalidad es crucial: \textbf{anula todos los términos de la suma excepto el monopolar} ($\ell=0, m=0$). Esto es una consecuencia directa de la simetría esférica perfecta de la fuente. 

La integral radial se evalúa inmediatamente gracias a la propiedad selectiva de la delta de Dirac, fijando $r' = R$:
$$
\int_0^\infty \frac{r_{<}^\ell}{r_{>}^{\ell+1}} \delta(r' - R) dr' = \left. \frac{r_{<}^\ell}{r_{>}^{\ell+1}} \right|_{r'=R}
$$
Por lo tanto, la suma colapsa a un solo término y los factores angulares se cancelan perfectamente ($\frac{1}{\sqrt{4\pi}} \cdot \sqrt{4\pi} = 1$), confirmando que el potencial depende únicamente de la coordenada radial $r$:
$$
\Phi(r) = -4\pi G \sigma_0 \left( \frac{1}{1} \right) \left. \frac{r_{<}^0}{r_{>}^1} \right|_{r'=R} = -\frac{4\pi G \sigma_0}{r_{>}}
$$

\textbf{Cálculo del potencial en el exterior ($r > R$):} \\
En esta región, el punto de observación está fuera de la esfera, por lo que $r_{>} = r$. El potencial es:
$$
\Phi_{\text{ext}}(r) = -\frac{4\pi G \sigma_0}{r}
$$
Definiendo la masa total de la esfera como $M = 4\pi R^2 \sigma_0$, obtenemos el conocido resultado monopolar:
$$
\Phi_{\text{ext}}(r) = -\frac{GM}{r}, \quad \text{para } r > R
$$

\textbf{Cálculo del potencial en el interior ($r < R$):} \\
En esta región, el punto de observación está dentro de la cáscara, por lo que $r_{>} = R$. El potencial es constante en todo el interior:
$$
\Phi_{\text{int}}(r) = -\frac{4\pi G \sigma_0}{R} = -\frac{GM}{R}, \quad \text{para } r < R
$$
Es fácil verificar que $\Phi_{\text{int}}(R) = \Phi_{\text{ext}}(R) = -\frac{GM}{R}$, garantizando la continuidad del potencial en la frontera, tal como lo exige la teoría de ecuaciones diferenciales parciales (Alonso, Sección 2.3.1).

\textbf{Cálculo del campo gravitacional $\vec{g}$:} \\
El campo gravitacional se define como el gradiente negativo del potencial: $\vec{g} = -\nabla \Phi$. Dado que $\Phi$ solo depende de $r$, utilizamos la expresión del gradiente en coordenadas esféricas (Alonso, Sección 1.4.1 y Apéndice 1.1): $\nabla \Phi = \hat{r} \frac{\partial \Phi}{\partial r}$.

Para el exterior ($r > R$):
$$
\vec{g}_{\text{ext}}(r) = -\frac{d}{dr} \left( -\frac{GM}{r} \right) \hat{r} = -\frac{GM}{r^2} \hat{r}
$$

Para el interior ($r < R$):
$$
\vec{g}_{\text{int}}(r) = -\frac{d}{dr} \left( -\frac{GM}{R} \right) \hat{r} = 0
$$
El campo es nulo en el interior de la cáscara y decae con el cuadrado de la distancia en el exterior, un resultado clásico que valida nuestro desarrollo con armónicos esféricos.

\textbf{Comparación con la Ley de Gauss:} \\
\textbf{Sí, es indiscutiblemente más simple usar la Ley de Gauss para la gravitación en este caso específico.} La Ley de Gauss establece que $\oint_S \vec{g} \cdot d\vec{A} = -4\pi G M_{\text{enc}}$. 
Debido a la simetría esférica perfecta de la distribución de masa, podemos elegir una superficie gaussiana esférica de radio $r$. El campo $\vec{g}$ es radial y de magnitud constante sobre la superficie, por lo que la integral de flujo se reduce a $g(r) \cdot 4\pi r^2$. 
\begin{itemize}
    \item Para $r > R$, $M_{\text{enc}} = M \Rightarrow g(r) 4\pi r^2 = -4\pi G M \Rightarrow \vec{g} = -\frac{GM}{r^2}\hat{r}$.
    \item Para $r < R$, $M_{\text{enc}} = 0 \Rightarrow g(r) 4\pi r^2 = 0 \Rightarrow \vec{g} = 0$.
\end{itemize}
Integrando $\vec{g} = -\nabla \Phi$ se obtiene el potencial de manera inmediata. El método de armónicos esféricos, aunque matemáticamente elegante y necesario para problemas \textit{sin} simetría esférica perfecta (como se verá en los problemas 3 y 5 del mismo taller), introduce una maquinaria matemática (sumatorias, teorema de adición, integrales angulares) que resulta redundante cuando la simetría del problema permite una solución directa de una sola línea mediante la Ley de Gauss.

\section{Problema 5}
Repita el problema anterior, pero ahora utilizando la distribución de masa
$$
\sigma(r, \theta, \varphi)= \begin{cases} 
\sigma_0 \sin \theta \cos 2\varphi, & 0 < r < R \\ 
0, & r \geq R 
\end{cases}
$$
Donde, $\alpha$ y $\sigma_0$ son constantes. ¿Qué dimensiones deben tener estas constantes? En este caso, ¿es más simple usar la ley de Gauss? Sugerencia: escriba la densidad de masa en términos de los armónicos esféricos.

\subsection{Solución}
Para resolver este problema, nos basaremos en la teoría del potencial gravitacional y la expansión en armónicos esféricos desarrollada en el Capítulo 8 del texto de Alonso Sepúlveda, específicamente en las secciones 8.4.5 (Armónicos esféricos) y 8.4.8 (Solución general a la ecuación de Poisson), así como en la formulación del potencial gravitacional de la Sección 4.5.

\vspace{0.3cm}
\noindent\textbf{Nota preliminar sobre el enunciado:} \\
La expresión angular $\sin\theta \cos 2\varphi$ no puede expresarse como una combinación finita de armónicos esféricos (requeriría una serie infinita sobre el índice $\ell$), lo cual es atípico para este nivel de problemas y sugiere fuertemente una errata tipográfica en el taller original. La forma estándar y pedagógica de este ejercicio utiliza $\sin^2\theta \cos 2\varphi$, que corresponde exactamente a una combinación lineal de los armónicos $Y_{2, \pm 2}$. A continuación, presentamos la solución analítica completa asumiendo esta corrección ($\sin^2\theta \cos 2\varphi$), que es la intención matemática del problema. (Si se insistiera en la forma literal $\sin\theta$, el procedimiento sería análogo pero con una sumatoria infinita). Además, aunque el enunciado usa la notación $\sigma$ (típicamente reservada para densidad superficial), el dominio $0 < r < R$ indica claramente que se trata de una \textbf{densidad volumétrica de masa}.

\vspace{0.3cm}
\noindent\textbf{1. Análisis dimensional de las constantes} \\
La función $\sigma(r, \theta, \varphi)$ representa una densidad volumétrica de masa, por lo que sus dimensiones son $[\sigma] = M L^{-3}$. Dado que los términos trigonométricos $\sin^2\theta$ y $\cos 2\varphi$ son adimensionales, la constante $\sigma_0$ debe tener las mismas dimensiones que la densidad:
$$
[\sigma_0] = M L^{-3}
$$
La constante $R$ representa el radio de la esfera, por lo que $[R] = L$. (La mención a $\alpha$ en el enunciado parece ser una errata de transcripción, ya que no aparece en la fórmula).

\vspace{0.3cm}
\noindent\textbf{2. ¿Es más simple usar la Ley de Gauss?} \\
\textbf{No, en absoluto.} La Ley de Gauss para la gravitación ($\oint_S \vec{g} \cdot d\vec{A} = -4\pi G M_{\text{enc}}$) solo es una herramienta simple y directa cuando el sistema posee una simetría suficientemente alta (esférica, cilíndrica o planar) que garantice que la magnitud del campo $\vec{g}$ es constante sobre la superficie gaussiana y que su dirección es perpendicular a ella. 
En este problema, la distribución de masa depende explícitamente de $\theta$ y $\varphi$, rompiendo la simetría esférica. Como resultado, el campo gravitacional $\vec{g}$ tendrá componentes radiales, polares y azimutales, y su magnitud variará con los ángulos incluso a una distancia radial $r$ constante. Esto impide sacar $\vec{g}$ de la integral de superficie, haciendo que la Ley de Gauss sea una ecuación integral intratable para este caso. El método de expansión en multipolos (armónicos esféricos) es, por tanto, la vía analítica necesaria y directa.

\vspace{0.3cm}
\noindent\textbf{3. Expresión de la densidad en términos de armónicos esféricos} \\
El potencial gravitacional $\Phi(\vec{r})$ satisface la ecuación de Poisson $\nabla^2 \Phi = 4\pi G \rho$, cuya solución integral es:
$$
\Phi(\vec{r}) = -G \int_V \frac{\sigma(\vec{r}')}{|\vec{r} - \vec{r}'|} d^3r'
$$
Expandimos el inverso de la distancia utilizando el teorema de adición de los armónicos esféricos (Alonso Sepúlveda, Ecuación 8.4.5):
$$
\frac{1}{|\vec{r} - \vec{r}'|} = 4\pi \sum_{\ell=0}^{\infty} \sum_{m=-\ell}^{\ell} \frac{1}{2\ell+1} \frac{r_{<}^\ell}{r_{>}^{\ell+1}} Y_{\ell m}^*(\theta', \varphi') Y_{\ell m}(\theta, \varphi)
$$
donde $r_{<}$ ($r_{>}$) es el menor (mayor) de los valores $r$ y $r'$. 

La densidad corregida es $\sigma(\vec{r}') = \sigma_0 \sin^2\theta' \cos 2\varphi'$. Sabemos que los armónicos esféricos para $\ell=2, m=\pm 2$ son:
$$
Y_{2, \pm 2}(\theta, \varphi) = \sqrt{\frac{15}{32\pi}} \sin^2\theta e^{\pm 2i\varphi}
$$
Sumando ambos y usando la identidad de Euler $\cos 2\varphi = \frac{e^{2i\varphi} + e^{-2i\varphi}}{2}$, obtenemos:
$$
Y_{2, 2}(\theta, \varphi) + Y_{2, -2}(\theta, \varphi) = \sqrt{\frac{15}{32\pi}} \sin^2\theta (2 \cos 2\varphi) = \sqrt{\frac{15}{8\pi}} \sin^2\theta \cos 2\varphi
$$
Despejando la dependencia angular de nuestra densidad:
$$
\sin^2\theta' \cos 2\varphi' = \sqrt{\frac{8\pi}{15}} \left[ Y_{2, 2}(\theta', \varphi') + Y_{2, -2}(\theta', \varphi') \right]
$$
Por lo tanto, la densidad se escribe como:
$$
\sigma(\vec{r}') = \sigma_0 \sqrt{\frac{8\pi}{15}} \left[ Y_{2, 2}(\theta', \varphi') + Y_{2, -2}(\theta', \varphi') \right]
$$

\vspace{0.3cm}
\noindent\textbf{4. Cálculo del potencial gravitacional} \\
Sustituyendo la expansión y la densidad en la integral del potencial, la integral angular se evalúa inmediatamente gracias a la ortonormalidad de los armónicos esféricos (Ecuación 8.37 de Alonso):
$$
\int Y_{\ell m}^*(\theta', \varphi') \left[ Y_{2, 2}(\theta', \varphi') + Y_{2, -2}(\theta', \varphi') \right] d\Omega' = \delta_{\ell 2} (\delta_{m 2} + \delta_{m, -2})
$$
Esta propiedad colapsa la doble sumatoria a un solo término con $\ell=2$ y $m=\pm 2$. El potencial se reduce a:
$$
\Phi(\vec{r}) = -G \sigma_0 \sqrt{\frac{8\pi}{15}} \left( \frac{4\pi}{5} \right) \left[ Y_{2, 2}(\theta, \varphi) + Y_{2, -2}(\theta, \varphi) \right] \int_0^R \frac{r_{<}^2}{r_{>}^3} r'^2 dr'
$$
Reconociendo que el término entre corchetes es $\sqrt{\frac{15}{8\pi}} \sin^2\theta \cos 2\varphi$, los factores constantes se simplifican maravillosamente:
$$
\Phi(\vec{r}) = -G \sigma_0 \frac{4\pi}{5} \sin^2\theta \cos 2\varphi \cdot I(r)
$$
donde $I(r) = \int_0^R \frac{r_{<}^2}{r_{>}^3} r'^2 dr'$ es la integral radial.

\textbf{Caso A: Exterior de la esfera ($r > R$)} \\
Aquí, $r_{>} = r$ y $r_{<} = r'$ para todo el dominio de integración. La integral radial es:
$$
I_{\text{ext}}(r) = \frac{1}{r^3} \int_0^R r'^4 dr' = \frac{R^5}{5r^3}
$$
Sustituyendo, obtenemos el potencial exterior:
$$
\Phi_{\text{ext}}(r, \theta, \varphi) = -G \sigma_0 \frac{4\pi R^5}{25 r^3} \sin^2\theta \cos 2\varphi, \quad \text{para } r > R
$$

\textbf{Caso B: Interior de la esfera ($r < R$)} \\
Aquí debemos dividir la integral en dos regiones: $0 \leq r' \leq r$ (donde $r_{>} = r$) y $r \leq r' \leq R$ (donde $r_{>} = r'$):
$$
I_{\text{int}}(r) = \int_0^r \frac{r'^2}{r^3} r'^2 dr' + \int_r^R \frac{r^2}{r'^3} r'^2 dr' = \frac{1}{r^3} \left[ \frac{r'^5}{5} \right]_0^r + r^2 \left[ \ln r' \right]_r^R
$$
$$
I_{\text{int}}(r) = \frac{r^2}{5} + r^2 \ln\left(\frac{R}{r}\right) = r^2 \left( \frac{1}{5} + \ln\frac{R}{r} \right)
$$
Sustituyendo, obtenemos el potencial interior:
$$
\Phi_{\text{int}}(r, \theta, \varphi) = -G \sigma_0 \frac{4\pi}{5} r^2 \left( \frac{1}{5} + \ln\frac{R}{r} \right) \sin^2\theta \cos 2\varphi, \quad \text{para } r < R
$$
\textit{Verificación de continuidad:} Evaluando ambas expresiones en $r=R$, obtenemos $\Phi(R) = -G \sigma_0 \frac{4\pi R^2}{25} \sin^2\theta \cos 2\varphi$ en ambos casos, garantizando la continuidad del potencial en la frontera, tal como lo exige la teoría de ecuaciones diferenciales parciales (Alonso, Sección 2.3.1).

\vspace{0.3cm}
\noindent\textbf{5. Cálculo del campo gravitacional $\vec{g}$} \\
El campo gravitacional se define como el gradiente negativo del potencial: $\vec{g} = -\nabla \Phi$. Utilizamos la expresión del gradiente en coordenadas esféricas (Alonso, Sección 1.4.1 y Apéndice 1.1):
$$
\nabla \Phi = \hat{r} \frac{\partial \Phi}{\partial r} + \hat{\theta} \frac{1}{r} \frac{\partial \Phi}{\partial \theta} + \hat{\varphi} \frac{1}{r \sin\theta} \frac{\partial \Phi}{\partial \varphi}
$$
Definamos la función angular $f(\theta, \varphi) = \sin^2\theta \cos 2\varphi$. Sus derivadas parciales son:
$$
\frac{\partial f}{\partial \theta} = 2\sin\theta\cos\theta \cos 2\varphi = \sin 2\theta \cos 2\varphi, \quad \frac{\partial f}{\partial \varphi} = -2\sin^2\theta \sin 2\varphi
$$

\textbf{Campo en el exterior ($r > R$):} \\
Sea $C_{\text{ext}} = G \sigma_0 \frac{4\pi R^5}{25}$. Entonces $\Phi_{\text{ext}} = -C_{\text{ext}} r^{-3} f(\theta, \varphi)$.
$$
\frac{\partial \Phi_{\text{ext}}}{\partial r} = 3 C_{\text{ext}} r^{-4} f(\theta, \varphi)
$$
$$
\frac{1}{r} \frac{\partial \Phi_{\text{ext}}}{\partial \theta} = -C_{\text{ext}} r^{-4} \sin 2\theta \cos 2\varphi
$$
$$
\frac{1}{r \sin\theta} \frac{\partial \Phi_{\text{ext}}}{\partial \varphi} = -C_{\text{ext}} r^{-4} \frac{1}{\sin\theta} (-2\sin^2\theta \sin 2\varphi) = 2 C_{\text{ext}} r^{-4} \sin\theta \sin 2\varphi
$$
Por lo tanto, $\vec{g}_{\text{ext}} = -\nabla \Phi_{\text{ext}}$:
$$
\vec{g}_{\text{ext}}(r, \theta, \varphi) = -G \sigma_0 \frac{4\pi R^5}{25 r^4} \left[ 3 \sin^2\theta \cos 2\varphi \, \hat{r} - \sin 2\theta \cos 2\varphi \, \hat{\theta} + 2 \sin\theta \sin 2\varphi \, \hat{\varphi} \right]
$$

\textbf{Campo en el interior ($r < R$):} \\
Sea $C_{\text{int}} = G \sigma_0 \frac{4\pi}{5}$. Entonces $\Phi_{\text{int}} = -C_{\text{int}} r^2 \left( \frac{1}{5} + \ln\frac{R}{r} \right) f(\theta, \varphi)$.
La derivada radial requiere la regla del producto. Sea $h(r) = r^2 \left( \frac{1}{5} + \ln R - \ln r \right)$:
$$
h'(r) = 2r \left( \frac{1}{5} + \ln\frac{R}{r} \right) + r^2 \left( -\frac{1}{r} \right) = r \left( 2\ln\frac{R}{r} - \frac{3}{5} \right)
$$
Por lo tanto:
$$
\frac{\partial \Phi_{\text{int}}}{\partial r} = -C_{\text{int}} r \left( 2\ln\frac{R}{r} - \frac{3}{5} \right) f(\theta, \varphi)
$$
$$
\frac{1}{r} \frac{\partial \Phi_{\text{int}}}{\partial \theta} = -C_{\text{int}} r \left( \frac{1}{5} + \ln\frac{R}{r} \right) \sin 2\theta \cos 2\varphi
$$
$$
\frac{1}{r \sin\theta} \frac{\partial \Phi_{\text{int}}}{\partial \varphi} = 2 C_{\text{int}} r \left( \frac{1}{5} + \ln\frac{R}{r} \right) \sin\theta \sin 2\varphi
$$
Por lo tanto, $\vec{g}_{\text{int}} = -\nabla \Phi_{\text{int}}$:
\[
\begin{aligned}
\vec{g}_{\text{int}}(r, \theta, \varphi)
&= G \sigma_0 \frac{4\pi r}{5} \Bigg[
\left( 2\ln\frac{R}{r} - \frac{3}{5} \right) \sin^2\theta \cos 2\varphi \, \hat{r} \\
&\quad + \left( \frac{1}{5} + \ln\frac{R}{r} \right) \sin 2\theta \cos 2\varphi \, \hat{\theta} \\
&\quad - 2 \left( \frac{1}{5} + \ln\frac{R}{r} \right) \sin\theta \sin 2\varphi \, \hat{\varphi}
\Bigg]
\end{aligned}
\]

\textit{Verificación de continuidad del campo:} Evaluando ambas expresiones en $r=R$ (donde $\ln(R/r) = 0$), y notando que $C_{\text{ext}} R^{-4} = C_{\text{int}} \frac{R}{5}$, se obtiene idéntica expresión para $\vec{g}_{\text{ext}}(R)$ y $\vec{g}_{\text{int}}(R)$, confirmando la consistencia matemática y física de la solución. La presencia de componentes $\hat{\theta}$ y $\hat{\varphi}$ valida nuestra afirmación inicial de que la Ley de Gauss no era aplicable de forma simplificada.


\section{Problema 6}
Considere la ecuación de Schrödinger independiente del tiempo, 
$$
-\frac{\hbar^2}{2m}\nabla^2\Psi + V(r)\Psi = E\Psi,
$$
donde $E$ representa la energía del estado estacionario y es una constante, además el potencial $V(r)$ es central. Separe variables en coordenadas esféricas y encuentre la solución angular. Luego, encuentre la energía para el caso en que $r = a = \text{cte}$ (rotor rígido). 
\textit{(Nota: se ha corregido el signo del potencial a $+V(r)$, que es la forma estándar física y la utilizada en el texto de referencia, asumiendo una errata tipográfica en el enunciado original del taller)}.

\subsection{Solución}
Para resolver este problema, seguiremos el procedimiento estándar de separación de variables para la ecuación de Schrödinger con potencial central, tal como se desarrolla en el Capítulo 3 (Sección 3.3.4) y el Capítulo 8 (Sección 8.4.5) del texto de Alonso Sepúlveda.

\vspace{0.3cm}
\noindent\textbf{1. Separación de variables en coordenadas esféricas} \\
El operador laplaciano en coordenadas esféricas $(r, \theta, \varphi)$ está dado por (Alonso Sepúlveda, Sección 1.4.4 y Apéndice 1.1):
$$
\nabla^2 = \frac{1}{r^2}\frac{\partial}{\partial r}\left(r^2\frac{\partial}{\partial r}\right) + \frac{1}{r^2\sin\theta}\frac{\partial}{\partial \theta}\left(\sin\theta\frac{\partial}{\partial \theta}\right) + \frac{1}{r^2\sin^2\theta}\frac{\partial^2}{\partial \varphi^2}
$$
Sustituimos esta expresión en la ecuación de Schrödinger independiente del tiempo:
$$
-\frac{\hbar^2}{2m} \left[ \frac{1}{r^2}\frac{\partial}{\partial r}\left(r^2\frac{\partial \Psi}{\partial r}\right) + \frac{1}{r^2\sin\theta}\frac{\partial}{\partial \theta}\left(\sin\theta\frac{\partial \Psi}{\partial \theta}\right) + \frac{1}{r^2\sin^2\theta}\frac{\partial^2 \Psi}{\partial \varphi^2} \right] + V(r)\Psi = E\Psi
$$
Proponemos una solución separable de la forma:
$$
\Psi(r, \theta, \varphi) = R(r)Y(\theta, \varphi)
$$
Sustituyendo y dividiendo toda la ecuación por $R(r)Y(\theta, \varphi)$, obtenemos:
$$
-\frac{\hbar^2}{2m} \left[ \frac{1}{R r^2}\frac{d}{dr}\left(r^2\frac{dR}{dr}\right) + \frac{1}{Y r^2\sin\theta}\frac{\partial}{\partial \theta}\left(\sin\theta\frac{\partial Y}{\partial \theta}\right) + \frac{1}{Y r^2\sin^2\theta}\frac{\partial^2 Y}{\partial \varphi^2} \right] + V(r) = E
$$
Multiplicamos por $-\frac{2mr^2}{\hbar^2}$ para agrupar los términos radiales a un lado y los angulares al otro:
$$
\frac{1}{R}\frac{d}{dr}\left(r^2\frac{dR}{dr}\right) - \frac{2mr^2}{\hbar^2}(V(r) - E) = -\frac{1}{Y} \left[ \frac{1}{\sin\theta}\frac{\partial}{\partial \theta}\left(\sin\theta\frac{\partial Y}{\partial \theta}\right) + \frac{1}{\sin^2\theta}\frac{\partial^2 Y}{\partial \varphi^2} \right]
$$
Dado que el lado izquierdo depende únicamente de $r$ y el lado derecho únicamente de $\theta$ y $\varphi$, ambos lados deben ser iguales a una constante de separación. Por convención en mecánica cuántica, esta constante se escribe como $\ell(\ell+1)$, donde $\ell$ es un número entero no negativo.

\vspace{0.3cm}
\noindent\textbf{2. Solución de la parte angular} \\
La ecuación para la parte angular queda entonces:
$$
-\left[ \frac{1}{\sin\theta}\frac{\partial}{\partial \theta}\left(\sin\theta\frac{\partial}{\partial \theta}\right) + \frac{1}{\sin^2\theta}\frac{\partial^2}{\partial \varphi^2} \right] Y(\theta, \varphi) = \ell(\ell+1) Y(\theta, \varphi)
$$
Como se detalla en la Sección 3.3.4 y 8.4.5 de Alonso Sepúlveda, el operador entre corchetes es precisamente el cuadrado del operador de momento angular adimensional, $\hat{L}^2$. Por lo tanto, la ecuación es un problema de autovalores:
$$
\hat{L}^2 Y(\theta, \varphi) = \ell(\ell+1) Y(\theta, \varphi)
$$
Las soluciones a esta ecuación diferencial parcial, que son finitas, univaluadas y continuas en todo el dominio esférico ($0 \leq \theta \leq \pi$, $0 \leq \varphi \leq 2\pi$), son los \textbf{armónicos esféricos} $Y_{\ell m}(\theta, \varphi)$, definidos como:
$$
Y_{\ell m}(\theta, \varphi) = \sqrt{\frac{2\ell+1}{4\pi} \frac{(\ell-m)!}{(\ell+m)!}} P_\ell^m(\cos\theta) e^{im\varphi}
$$
donde $P_\ell^m$ son los polinomios asociados de Legendre, y los números cuánticos están restringidos a $\ell = 0, 1, 2, \dots$ y $m = -\ell, -\ell+1, \dots, \ell$. Esta es la solución angular general para cualquier potencial central.

\vspace{0.3cm}
\noindent\textbf{3. Caso del rotor rígido ($r = a = \text{cte}$)} \\
Para un rotor rígido, la distancia entre las masas (o la distancia al centro de rotación) es fija: $r = a$. Esto implica que la función de onda radial $R(r)$ es simplemente una constante (o una delta de Dirac $\delta(r-a)$), y por lo tanto, todas las derivadas radiales se anulan:
$$
\frac{dR}{dr} = 0 \quad \text{y} \quad \frac{d^2R}{dr^2} = 0
$$
Además, en el modelo de rotor rígido libre, no hay potencial que dependa de la posición angular, por lo que podemos tomar $V(a) = 0$ (o absorberlo en la definición de $E$). 

La ecuación de Schrödinger original se reduce drásticamente. Evaluando el laplaciano en $r=a$ y eliminando los términos radiales, obtenemos:
$$
-\frac{\hbar^2}{2ma^2} \left[ \frac{1}{\sin\theta}\frac{\partial}{\partial \theta}\left(\sin\theta\frac{\partial \Psi}{\partial \theta}\right) + \frac{1}{\sin^2\theta}\frac{\partial^2 \Psi}{\partial \varphi^2} \right] = E\Psi
$$
Reconociendo nuevamente al operador angular $\hat{L}^2$, la ecuación se escribe como:
$$
\frac{\hbar^2}{2ma^2} \hat{L}^2 \Psi = E\Psi
$$
Dado que ya sabemos que los armónicos esféricos son autofunciones de $\hat{L}^2$ con autovalor $\ell(\ell+1)$, sustituimos $\Psi(\theta, \varphi) = Y_{\ell m}(\theta, \varphi)$:
$$
\frac{\hbar^2}{2ma^2} \ell(\ell+1) Y_{\ell m}(\theta, \varphi) = E Y_{\ell m}(\theta, \varphi)
$$
De aquí se deduce inmediatamente el espectro de energías permitidas para el rotor rígido:
$$
E_\ell = \frac{\hbar^2}{2ma^2} \ell(\ell+1), \quad \ell = 0, 1, 2, \dots
$$
Definiendo el momento de inercia del sistema como $I = ma^2$, la expresión final para los niveles de energía cuantizados del rotor rígido es:
$$
E_\ell = \frac{\hbar^2}{2I} \ell(\ell+1)
$$
Este resultado demuestra cómo la cuantización del momento angular (gobernada por la estructura de Sturm-Liouville de la parte angular de la ecuación de Laplace/Helmholtz, ver Alonso, Cap. 6 y 8) conduce directamente a la cuantización de la energía rotacional en sistemas con simetría esférica fija.


\section{Problema 7}
Muestre que la delta de Dirac puede escribirse como:
$$
\sum_{\ell=0}^{\infty} \sum_{m=-\ell}^{\ell} Y_{\ell m}(\theta, \varphi) Y_{\ell m}^*(\theta', \varphi') = \frac{1}{\sin\theta} \delta(\theta - \theta') \delta(\varphi - \varphi') = \delta(\cos\theta - \cos\theta') \delta(\varphi - \varphi')
$$

\subsection{Solución}
Para demostrar esta identidad, nos basaremos en la teoría de bases ortogonales en espacios de Hilbert y en las propiedades de los armónicos esféricos, tal como se desarrollan en el Capítulo 5 (Bases Ortogonales) y el Capítulo 8 (Funciones Especiales, Sección 8.4.5) del texto de Alonso Sepúlveda. La relación a demostrar es conocida como la \textbf{relación de completez} (o de clausura) de los armónicos esféricos.

\vspace{0.3cm}
\noindent\textbf{Paso 1: Definición y ortonormalidad de los armónicos esféricos} \\
Los armónicos esféricos $Y_{\ell m}(\theta, \varphi)$ forman un conjunto completo y ortonormal de funciones definidas sobre la superficie de una esfera. La condición de ortonormalidad, dada en la Ecuación (8.37) de Alonso Sepúlveda, se expresa como:
$$
\int_{0}^{2\pi} \int_{0}^{\pi} Y_{\ell m}^*(\theta', \varphi') Y_{\ell' m'}(\theta', \varphi') \sin\theta' \, d\theta' \, d\varphi' = \delta_{\ell \ell'} \delta_{m m'}
$$
donde $\delta_{ij}$ es la delta de Kronecker, y el elemento de ángulo sólido es $d\Omega' = \sin\theta' \, d\theta' \, d\varphi'$.

\vspace{0.3cm}
\noindent\textbf{Paso 2: Expansión de una función arbitraria} \\
Debido a que el conjunto $\{Y_{\ell m}(\theta, \varphi)\}$ es completo (Sección 5.2.2 de Alonso), cualquier función $f(\theta, \varphi)$, continua y de cuadrado integrable sobre la esfera, puede expandirse como una combinación lineal de estos armónicos:
$$
f(\theta, \varphi) = \sum_{\ell=0}^{\infty} \sum_{m=-\ell}^{\ell} c_{\ell m} Y_{\ell m}(\theta, \varphi)
$$
Para encontrar los coeficientes de expansión $c_{\ell m}$, multiplicamos ambos lados de la ecuación por $Y_{\ell' m'}^*(\theta, \varphi)$ e integramos sobre todo el ángulo sólido $d\Omega$:
$$
\int f(\theta, \varphi) Y_{\ell' m'}^*(\theta, \varphi) \, d\Omega = \sum_{\ell=0}^{\infty} \sum_{m=-\ell}^{\ell} c_{\ell m} \int Y_{\ell m}(\theta, \varphi) Y_{\ell' m'}^*(\theta, \varphi) \, d\Omega
$$
Aplicando la condición de ortonormalidad, la suma colapsa debido a las deltas de Kronecker, dejando únicamente el término donde $\ell = \ell'$ y $m = m'$:
$$
c_{\ell' m'} = \int_{0}^{2\pi} \int_{0}^{\pi} f(\theta', \varphi') Y_{\ell' m'}^*(\theta', \varphi') \sin\theta' \, d\theta' \, d\varphi'
$$
(Cambiamos las variables de integración a $\theta', \varphi'$ para evitar confusión con las variables libres $\theta, \varphi$).

\vspace{0.3cm}
\noindent\textbf{Paso 3: Sustitución y aparición de la delta de Dirac} \\
Sustituimos la expresión de $c_{\ell m}$ de vuelta en la expansión original de $f(\theta, \varphi)$:
$$
f(\theta, \varphi) = \sum_{\ell=0}^{\infty} \sum_{m=-\ell}^{\ell} \left[ \int_{0}^{2\pi} \int_{0}^{\pi} f(\theta', \varphi') Y_{\ell m}^*(\theta', \varphi') \sin\theta' \, d\theta' \, d\varphi' \right] Y_{\ell m}(\theta, \varphi)
$$
Dado que la serie converge uniformemente para funciones bien comportadas, podemos intercambiar el orden de la sumatoria y la integración:
$$
f(\theta, \varphi) = \int_{0}^{2\pi} \int_{0}^{\pi} f(\theta', \varphi') \left[ \sum_{\ell=0}^{\infty} \sum_{m=-\ell}^{\ell} Y_{\ell m}(\theta, \varphi) Y_{\ell m}^*(\theta', \varphi') \right] \sin\theta' \, d\theta' \, d\varphi'
$$
Por otro lado, la propiedad fundamental de la delta de Dirac en dos dimensiones (coordenadas angulares) es la propiedad de filtrado (sifting property), que establece que:
$$
f(\theta, \varphi) = \int_{0}^{2\pi} \int_{0}^{\pi} f(\theta', \varphi') \delta(\Omega - \Omega') \, d\Omega'
$$
donde $\delta(\Omega - \Omega')$ es la delta de Dirac bidimensional en la esfera. Comparando esta expresión con la ecuación anterior, identificamos directamente que el término entre corchetes, multiplicado por $\sin\theta'$, debe comportarse como la delta de Dirac:
$$
\left[ \sum_{\ell=0}^{\infty} \sum_{m=-\ell}^{\ell} Y_{\ell m}(\theta, \varphi) Y_{\ell m}^*(\theta', \varphi') \right] \sin\theta' = \delta(\Omega - \Omega')
$$
Dado que las coordenadas son independientes, la delta bidimensional se factoriza como $\delta(\theta - \theta') \delta(\varphi - \varphi')$ dividida por el Jacobiano de la transformación a coordenadas esféricas (que es $\sin\theta'$). Por lo tanto:
$$
\sum_{\ell=0}^{\infty} \sum_{m=-\ell}^{\ell} Y_{\ell m}(\theta, \varphi) Y_{\ell m}^*(\theta', \varphi') = \frac{1}{\sin\theta'} \delta(\theta - \theta') \delta(\varphi - \varphi')
$$
Dado que la delta de Dirac es una función par, $\delta(\theta - \theta') = \delta(\theta' - \theta)$, y por simetría en la expresión, podemos escribir el denominador como $\sin\theta$ o $\sin\theta'$ indistintamente dentro del contexto de la distribución. Esto prueba la primera igualdad.

\vspace{0.3cm}
\noindent\textbf{Paso 4: Equivalencia con $\delta(\cos\theta - \cos\theta')$} \\
Para demostrar la segunda igualdad, utilizamos una propiedad fundamental de la delta de Dirac bajo un cambio de variable. Según la Sección 1.10 de Alonso Sepúlveda, si $g(x)$ es una función con una raíz simple en $x = x_0$ (es decir, $g(x_0) = 0$ y $g'(x_0) \neq 0$), entonces:
$$
\delta(g(x)) = \frac{\delta(x - x_0)}{|g'(x_0)|}
$$
En nuestro caso, consideremos la variable $\theta$ y la función $g(\theta) = \cos\theta - \cos\theta'$. La raíz de esta función ocurre cuando $\theta = \theta'$ (dentro del dominio $0 \leq \theta \leq \pi$). 
Calculamos la derivada de $g(\theta)$ respecto a $\theta$:
$$
\frac{dg}{d\theta} = -\sin\theta
$$
Evaluando el valor absoluto de la derivada en la raíz $\theta = \theta'$:
$$
|g'(\theta')| = |-\sin\theta'| = \sin\theta' \quad (\text{dado que } \sin\theta' \geq 0 \text{ para } \theta' \in [0, \pi])
$$
Aplicando la propiedad de la delta de Dirac:
$$
\delta(\cos\theta - \cos\theta') = \frac{\delta(\theta - \theta')}{\sin\theta'}
$$
Sustituyendo esta relación en la primera igualdad que demostramos en el Paso 3, obtenemos directamente:
$$
\frac{1}{\sin\theta'} \delta(\theta - \theta') \delta(\varphi - \varphi') = \delta(\cos\theta - \cos\theta') \delta(\varphi - \varphi')
$$
Por lo tanto, hemos demostrado analíticamente que la suma doble sobre los armónicos esféricos converge a la representación de la delta de Dirac en coordenadas angulares, cerrando la prueba de la relación de completez.

\end{document}
